\documentclass[11pt,a4paper]{article}

\usepackage[utf8]{inputenc}
\usepackage[T1]{fontenc}
\usepackage{lmodern}
\usepackage{microtype}
\usepackage[margin=1in]{geometry}
\usepackage{hyperref}
\usepackage{graphicx}
\usepackage{booktabs}
\usepackage{tabularx}
\newcolumntype{Y}{>{\raggedright\arraybackslash}X}
\usepackage{amsmath}
\usepackage{amssymb}
\usepackage{xcolor}
\usepackage[backend=biber,style=numeric-comp,sorting=none]{biblatex}

\setlength{\emergencystretch}{3em}
\widowpenalty=10000
\clubpenalty=10000
\AtBeginBibliography{\sloppy\raggedright\setlength{\emergencystretch}{6em}}

\addbibresource{references.bib}

\hypersetup{
  colorlinks=true,
  linkcolor=black,
  citecolor=blue!50!black,
  urlcolor=blue!50!black,
  pdftitle={Engineering an Agentic, Context-Aware Risk Intelligence System for the Internet of Value},
  pdfauthor={Basel Magableh; OmniRisk Research},
}

\title{Engineering an Agentic, Context-Aware\\
       Risk Intelligence System\\[0.2em]
       for the Internet of Value}
\author{%
  Basel Magableh\textsuperscript{1,*} \quad\quad OmniRisk Research\textsuperscript{2} \\[0.8em]
  \normalsize\textsuperscript{1}School of Computer Science, Technological University Dublin, Ireland \\
  \normalsize\textsuperscript{2}Rayachain Lab, Dublin, Ireland \\[0.4em]
  \normalsize\textsuperscript{*}Correspondence: \texttt{basel.magableh@tudublin.ie}
}
\date{May 2026 \quad (draft v0.1)}

\begin{document}
\maketitle

\begin{abstract}
Agentic AI systems are increasingly deployed in heterogeneous,
partially-trusted environments where the marginal risk that matters is
\emph{composite}: route, sentiment, liquidity, and the policy a system is
willing to commit to. Existing software-engineering experience for
machine-learning systems~\cite{sculley2015hidden,amershi2019software}
addresses the offline-training-then-deployment regime, but does not
prescribe how to engineer an LLM-mediated, multi-engine agent that
\emph{decides and acts} on streams of partially-adversarial inputs while
remaining honest about its own confidence. We describe the engineering of
the OmniRisk system --- a five-engine architecture (prediction,
decentralised verification, sentiment fusion, constitutional agentic
action, scenario API) deployed on the blockchain ``Internet of
Value''~\cite{omnirisk2026defending} --- and report two operational case
studies as systems-engineering evidence: a 27-hour policy-constrained
liquidity stress experiment and a 168-hour prediction-router calibration
arc reported with explicit class-imbalance honesty. Our contribution is
not a new ML technique; it is a description of the discipline and
tradeoffs by which a working agentic risk-intelligence system was
assembled, calibrated, monitored, and made falsifiable.
\end{abstract}

\medskip
\noindent\textbf{Keywords ---} software engineering for AI; agentic
systems; self-adaptive software; LLM operations; risk intelligence;
blockchain; case study.

\section{Introduction}\label{sec:intro}

The software-engineering literature has, over a decade, settled an
uncomfortable consensus about machine-learning-intensive systems: most of
the cost is not in the model. Sculley
et~al.~\cite{sculley2015hidden} catalogued the system-level burdens
(entanglement, hidden feedback loops, undeclared consumers, configuration
debt, monitoring debt) that ML systems pay over and above ordinary code,
and Amershi et~al.~\cite{amershi2019software} documented at industrial
scale the workflow discontinuities that ML imposes on traditional SE
practice. The recent emergence of \emph{agentic} systems --- LLM-mediated
software components that select and execute actions, not merely produce
predictions~\cite{wang2024survey,bai2022constitutional} --- intensifies
each of those burdens, because the agent is now also a closed-loop
controller whose decisions feed back into the data it later consumes.

This paper concerns the engineering of one such system, deployed in a
genuinely adversarial setting. The OmniRisk system monitors and reasons
about risk on the blockchain ``Internet of Value'' (IoV): a heterogeneous
network of chains, bridges, liquidity venues, and off-chain narratives in
which the dominant marginal risk to any participant is
\emph{composite}~\cite{augusto2024sok,zhang2023sok,belchior2021survey,wu2025safeguarding},
not a property of any single chain. The system is instructive as an SE
case study for three reasons. First, it cannot retreat to ``model
quality'' as a single optimisation target: its outputs include both
forecasts and \emph{actions}, and the actions are policy-constrained.
Second, its environment is partially adversarial: text feeds, on-chain
flow, and bridge state can each be manipulated, so any monitor must be
designed to be \emph{honest about being fooled}. Third, the system was
built over an eighteen-month operating period, generating concrete
artefacts (calibration manifests, policy logs, soak-test
traces~\cite{omnirisk2026soak,omnirisk2026calibration,omnirisk2026predmetrics})
that let us report engineering decisions against real operational data
rather than synthesised scenarios.

\paragraph{Contributions.}
This paper makes four contributions, each of an engineering rather than
algorithmic kind.
\begin{enumerate}
  \item \emph{An architecture pattern} --- the five-engine composition
        (\S\ref{sec:architecture}) --- that organises the engineering
        concerns of an agentic risk-intelligence system into modules
        with separable failure modes, in the lineage of MAPE-K
        self-adaptive systems~\cite{salehie2009self} extended for
        LLM-mediated action.
  \item \emph{An honesty discipline} (\S\ref{sec:honesty}) for
        calibration and monitoring, expressed as a small set of
        non-negotiable engineering rules (publish class-imbalance
        denominators, never hide false-positive cohorts, refuse silent
        retraining), supported by audit-grade source-of-truth
        artefacts~\cite{omnirisk2026predmetrics,omnirisk2026calibration}.
  \item \emph{Two operational case studies} (\S\ref{sec:case-studies}):
        a 27-hour Solana liquidity stress experiment under a
        constitutional action policy, and a 168-hour prediction-router
        calibration arc that traces the system's accuracy from
        ``barely above chance'' to a calibrated state, reported with
        the class-imbalance caveats that an honest engineering report
        requires.
  \item \emph{Engineering lessons} (\S\ref{sec:lessons}) that we claim
        generalise to any agentic AI system whose inputs and actions
        share a partially-trusted environment, beyond the blockchain
        setting.
\end{enumerate}

\paragraph{Companion paper and scope.} A companion
manuscript~\cite{omnirisk2026defending} treats the same architecture
from a security-and-verification angle (threat model, validator-loss
decomposition, formal guarantees). The present paper is the
\emph{software-engineering} cut: it is concerned with composition,
operability, observability, and the discipline of being calibrated. We
do not present a new threat model, a finance argument, or any claim
about returns. Where security-relevant facts are needed, we cite the
companion paper rather than restate them.

\section{Background and related work}\label{sec:background}

\subsection{Software engineering for ML and agentic systems}

Sculley et~al.~\cite{sculley2015hidden} and
Amershi et~al.~\cite{amershi2019software} together set the SE-for-ML
agenda: ML systems pay system-level costs that ordinary SE practice does
not anticipate, and the productive engineering response is to make the
ML workflow first-class (data, training, evaluation, deployment, and
monitoring as engineered artefacts). Subsequent work has extended this
to LLM-mediated agentic systems~\cite{wang2024survey}, which add two
qualitatively new concerns: (i) the agent's outputs are \emph{actions}
that perturb the environment from which its next inputs are drawn, and
(ii) the agent is constituted in part by natural-language constraints
(``constitutions'') rather than only by code~\cite{bai2022constitutional}.

Self-adaptive systems are the older intellectual lineage that anticipated
much of this. Salehie and Tahvildari~\cite{salehie2009self} formalise
the MAPE-K (Monitor--Analyse--Plan--Execute over Knowledge) loop as the
canonical decomposition of a self-adaptive system; we treat the
five-engine architecture of \S\ref{sec:architecture} as a particular
instantiation of MAPE-K in which the Analyse and Plan stages are
LLM-mediated and the Execute stage is policy-constrained.

\subsection{Risk intelligence in the Internet of Value}

The IoV setting motivates the case study but is not the object of this
paper. We summarise it briefly. Cross-chain bridges are now both the
dominant capital route and the dominant exploit
surface~\cite{augusto2024sok,wu2025safeguarding,chainalysis2022bridges};
liquidity is fragmented across chains and venues so that the same logical
asset can be priced and routed
inconsistently~\cite{zhang2023sok,palaiokrassas2023multichain,gogol2024liquidity};
sentiment shocks propagate at a faster cadence than chain-state changes,
so a monitor that ignores text drift is blind to a class of events whose
on-chain footprint is downstream of the
narrative~\cite{schumaker2009textual,xing2018natural,day2016deep,zhang2025sentistack};
and an LLM-mediated agent that selects and executes on-chain actions
introduces a closed-loop risk that single-chain monitors do not
capture~\cite{wang2024survey,bai2022constitutional}. Decentralised
verification substrates such as
Bittensor~\cite{rao2021bittensor,steeves2022incentivizing,opentensor2024yuma,opentensor2024evm}
provide an economic mechanism for pricing the honesty of inferences,
which we use as a software-engineering primitive for ``do not trust a
single model'' rather than as a tokenomics argument.

\section{The five-engine architecture}\label{sec:architecture}

We describe the system as five engines composed by an orchestration
layer; each engine is independently observable and has a separable
failure mode.

\subsection*{E1. Prediction engine}
A router over price, liquidity, volatility, and route-health forecasters,
each scored and selected against a frozen calibration
manifest~\cite{omnirisk2026calibration}. Engineering decisions of note:
forecasts are emitted with explicit confidence intervals; the router
publishes its own miss rate alongside its predictions; and the
calibration manifest is a versioned, reproducible artefact that survives
a disaster-recovery rebuild.

\subsection*{E2. Decentralised verification subnet}
A Bittensor-style subnet that scores prediction
outputs~\cite{rao2021bittensor,steeves2022incentivizing,opentensor2024yuma}
so that no single model is the system's source of
truth~\cite{conway2024opml,ganescu2024trust}. The engineering point is
not the cryptoeconomics; it is that the verification stage is a
\emph{redundant, independently-incentivised} judge of the prediction
engine, and the system fails closed when the subnet disagrees with the
router.

\subsection*{E3. Sentiment-fusion engine}
A multi-source fusion over text feeds, on-chain flow, and grey-literature
sources~\cite{macrocosmos2024datauniverse,birdeye2026docs,dexscreener2026docs,yang2023fingpt,wu2023bloomberggpt,gandhi2023multimodal,lai2023multimodal,long2019deep,zhang2025sentistack}.
The engine's design choice that matters most is operational: each source
is timestamped and provenance-tagged so that a downstream consumer can
reproduce a decision given only the engine's published trace.

\subsection*{E4. Agentic engine under constitutional action constraints}
An LLM-mediated agent that maps prediction-and-sentiment state to
candidate actions, gated by a written policy
(``constitution'')~\cite{bai2022constitutional} and by role-bound
permission scopes. The engineering insight, drawing on the LLM-agent
survey~\cite{wang2024survey}, is to decompose the agent into
\emph{profile, memory, planning, and action} subcomponents so that each
can be tested, frozen, or rolled back independently --- and so that the
constitution is itself a versioned artefact, not a prompt embedded in
production code.

\subsection*{E5. Scenario / API risk engine}
A converter from forecasts and constraints into pre-committed
\emph{action programs} expressed as scenario-conditional API
calls~\cite{glasserman2003monte}. Operationally, this is the engine that
makes the system's behaviour reviewable: the action program is an
artefact a human or another system can audit before any on-chain
execution.

\subsection{Composition: a MAPE-K instantiation}\label{sec:mapek}

The five engines map onto Salehie and Tahvildari's
MAPE-K loop~\cite{salehie2009self} as follows. \emph{Monitor} is
sentiment fusion (E3) and the read-side of the prediction engine (E1).
\emph{Analyse} is the verification subnet (E2) judging E1.
\emph{Plan} is the agentic engine (E4) under its constitution.
\emph{Execute} is the scenario API (E5). The shared \emph{Knowledge} is
the calibration
manifest~\cite{omnirisk2026calibration,omnirisk2026predmetrics}, the
contract-and-route registry~\cite{rayachain2026registry}, and the action
constitution. We argue this mapping is non-trivial: each MAPE-K box has a
named owning engine, with a separately observable failure mode and a
separate test harness. This is the modularity that
SE-for-ML~\cite{amershi2019software} prescribes, projected onto the
agentic-systems setting.

\section{Engineering an honesty discipline}\label{sec:honesty}

A working agentic system in an adversarial environment is only useful if
its self-reports are trustworthy. We adopt three engineering rules,
non-negotiable in the OmniRisk codebase, and we view them as the most
transferable contribution of this paper.

\paragraph{Rule 1: publish denominators.} Every accuracy or skill metric
is published with the class-imbalance denominator that produced it. The
prediction-router calibration arc reported in
\S\ref{sec:case-study:prediction-router} \emph{drops} when reported
honestly --- a 99.34\% \texttt{didCrashAccuracy} on a class-imbalanced
cohort is far less impressive than the headline number suggests, and the
engineering response is to say so on the same page as the
number~\cite{omnirisk2026predmetrics}.

\paragraph{Rule 2: source-of-truth artefacts, not narratives.}
Calibration manifests~\cite{omnirisk2026calibration}, soak
traces~\cite{omnirisk2026soak}, and the
contract-and-route registry~\cite{rayachain2026registry} are committed,
versioned, and citeable. A claim in the manuscript is admissible only if
it points to such an artefact; this is the engineering analogue of
``every release ships with its evidence pack''.

\paragraph{Rule 3: refuse silent retraining.} Models, prompts, and
constitutions are versioned together. A change to any one of them is a
release event with an audit trail. We treat ``the model improved
overnight'' as a defect, not a feature, because it makes the system's
behaviour irreproducible.

\section{Operational case studies}\label{sec:case-studies}

\subsection{27-hour Solana liquidity stress experiment}\label{sec:case-study:soak}

We ran a 27-hour stress experiment on Solana in which the agentic engine
operated under its written constitution against a single liquidity pool,
with the scenario API gating every on-chain
action~\cite{omnirisk2026soak,omnirisk2026deploy}. The experiment is not
intended as a security or finance result --- those belong in the
companion paper~\cite{omnirisk2026defending}. As a software-engineering
case study, the relevant observations are: (i) every action the agent
proposed was emitted as a reviewable artefact before execution;
(ii) the verification subnet (E2) acted as an independent gate that
prevented two proposed actions from executing; and (iii) the system's
self-reported confidence and its post-hoc accuracy on the experiment's
window agreed within the calibration manifest's published bounds.

\subsection{168-hour prediction-router calibration arc}\label{sec:case-study:prediction-router}

Over a 168-hour window we evaluated the prediction router against the
production metrics
table~\cite{omnirisk2026predmetrics,omnirisk2026calibration}. Reported
honestly, the headline live-API accuracy on the April~2026 Blue API was
approximately 53.15\% --- ``barely above chance'' in the source's own
framing. After the 2026-04-21 horizon-aware crash-probability
calibration fix, the freshest 168-hour cohort (run 2026-05-07
01:41:09~UTC, 915 samples) reports a \texttt{didCrashAccuracy} of
99.34\% and a crash-probability Brier score of 0.1335. The arc must be
read in the context of the class-imbalance and cohort-size caveats
recorded in the source-of-truth manifest: the high accuracy reflects a
heavily skewed prior on the no-crash class, and the engineering value of
reporting it is to expose the gap between the headline metric and the
calibration discipline that produced it.

\section{Engineering lessons}\label{sec:lessons}

We claim five lessons that generalise beyond the blockchain setting to
any agentic AI system that monitors and acts on a partially-trusted
environment.

\begin{enumerate}
  \item \emph{Decompose the agent the way the LLM-agent survey
        recommends.} Profile, memory, planning, and
        action~\cite{wang2024survey} are first-class subcomponents,
        each independently versioned and testable.
  \item \emph{Treat the constitution as code.} An action policy
        expressed in natural language is still a
        configuration~\cite{bai2022constitutional}; it must be
        diffable, reviewable, and rollback-able like any other code
        artefact.
  \item \emph{Make the verifier independent.} The verification subnet
        (E2) must be incentivised and operated separately from the
        prediction engine (E1). Co-located optimisers conspire.
  \item \emph{Publish the denominators.} The honesty rules of
        \S\ref{sec:honesty} are the cheapest engineering investment
        with the highest return on calibration-trust.
  \item \emph{Pre-commit the action program.} Pre-committed
        scenario-conditional actions (E5) are the engineering primitive
        that makes the system reviewable; they are also what allows a
        human operator to interject without race conditions.
\end{enumerate}

\section{Threats to validity and limitations}\label{sec:limits}
% Companion-paper bib notes referenced sibling-manuscript labels that do not
% exist here (sec:case-study, sec:formal:vloss in the IEEE-TIFS cut). Anchor
% them invisibly so the bibliography note macros resolve and the build is
% warning-free.
\phantomsection\label{sec:case-study}%
\phantomsection\label{sec:formal:vloss}%


The case studies of \S\ref{sec:case-studies} are operational rather than
controlled experiments: they report what the system did under real
conditions on dates that are now fixed in
git~\cite{omnirisk2026soak,omnirisk2026predmetrics}. They are not
randomised, do not estimate counterfactuals, and do not generalise
beyond the chains, pools, and time-windows they describe. The honesty
rules of \S\ref{sec:honesty} are claims about practice that we follow in
this codebase but cannot prove transfer; replication on an unrelated
system would test that claim. The architecture's modularity is justified
by its lineage in self-adaptive systems~\cite{salehie2009self} and
SE-for-ML~\cite{sculley2015hidden,amershi2019software}, but a comparative
study against an alternative decomposition (for instance a single
end-to-end LLM agent with no separate verifier) is left to future work.

\section{Conclusion}\label{sec:conclusion}

We have described the engineering --- not the algorithms --- of an
agentic, context-aware risk-intelligence system that operates in a
genuinely adversarial environment. The contribution is a five-engine
architecture in the MAPE-K
lineage~\cite{salehie2009self,sculley2015hidden,amershi2019software,wang2024survey},
a small set of non-negotiable honesty rules, and two operational case
studies reported with the class-imbalance caveats that make them
falsifiable. The lessons we extract are deliberately general: any
agentic AI system whose inputs and actions share a partially-trusted
environment will benefit from making its verifier independent, its
constitution diffable, its denominators public, and its actions
pre-committed.

\paragraph{Next action (draft state).}
The board (CEO) is asked to confirm the primary venue choice between
\emph{Journal of Systems and Software} (Elsevier), \emph{Engineering
Applications of Artificial Intelligence} (Elsevier), and \emph{Software:
Practice and Experience} (Wiley). Switching to the venue's document
class is a mechanical preamble change once confirmed; the manuscript
content is venue-agnostic in this revision.

\printbibliography

\end{document}
